\chapter{Introduction}
\label{ch:intro}

\section{Motivation}
\label{sec:intro:motivation}

Graphs arise across a wide range of scientific and technological domains. Molecules can be represented as atoms connected by chemical bonds; citation networks as documents linked by references; social networks as individuals connected through interactions; and biological systems as proteins, genes, or metabolites related through functional processes. More generally, many complex systems are naturally described as collections of entities together with the relationships among them. In these settings, the graph is not merely a convenient storage format: its connectivity is an essential part of the object being modelled.

This thesis studies three closely connected problems in graph machine learning:
how to construct graph representations that expose structural information, how
to generate new graphs under graph-level conditions, and how to evaluate graph
generative models with explicit reliability checks. Although these problems are
often considered separately, they influence one another. A generative model can
only be guided by information that has been represented effectively; generated
graphs can only be trusted when their quality and task-relevant utility can be
assessed; and such assessments depend on representations that expose the
structural properties of interest.

Graph data present challenges that distinguish them from vectors, images, and sequences. Graphs have no canonical ordering of their nodes, may vary substantially in size, and often contain irregular local structures~\citep{bronstein2017geometric}. Their defining properties may depend on small motifs, extended neighbourhoods, or long-range connectivity across distant regions. These characteristics motivate expressive learning methods such as graph neural networks and graph diffusion models, but they also complicate interpretability, robustness, and evaluation.

Graph kernels provide a complementary approach to learning from graph-struc\-tured data. Over the past two decades, numerous kernels have been developed to compare graphs through structural components such as walks, paths, neighbourhoods, subtrees, and graphlets~\citep{nikolentzos2021graph}. Their explicit and deterministic representations can support feature-level inspection and sample efficiency. However, many classical graph kernels were designed primarily for discrete node labels and therefore handle continuous attributes only through discretisation or other indirect mechanisms. Neural approaches can learn more flexible representations from attributed data, but their behaviour may be sensitive to optimisation and architectural choices and can be more difficult to inspect.

This thesis is motivated by the view that reliable graph machine learning should combine the strengths of these traditions. It develops explicit representations that retain meaningful graph structure, conditional generative models whose outputs can be guided by graph-level information, and evaluation protocols that distinguish genuine structural and task-relevant utility from superficial similarity.


\section{Problem Setting}
\label{sec:intro:problem}

The central object considered in this thesis is a graph
\[
    G=(V,E),
\]
possibly equipped with node labels, edge labels, or continuous node attributes.
The thesis focuses on graph-level and dataset-level questions rather than on a
single supervised prediction task. In particular, it considers:

\begin{enumerate}[label=(\roman*)]
    \item representation: constructing graph feature maps that preserve
    structural and attribute information while remaining computationally
    practical and interpretable;
    \item generation: producing labelled graphs from graph-level conditioning
    information using a staged denoising and decoding process; and
    \item evaluation: comparing generative models using diagnostic,
    statistically controlled measures rather than relying on a single heuristic
    score.
\end{enumerate}

These three components form the methodological arc of the thesis, but they
should not be read as a single mandatory end-to-end system. The contributions
are operationally independent: NSPPK is a graph representation method, GCDG is
a graph-conditioned generator, and RankGen is a general diagnostic evaluation
framework. They address successive decisions that arise in graph generative
modelling: what graph information is represented, how such information can
condition generation, and how generated data should be assessed.

\section{Aim, Objectives, and Research Questions}
\label{sec:intro:questions}

The overall aim of this thesis is to develop graph representation,
graph-conditioned generation, and diagnostic evaluation methods that are
explicit about structural information, conditioning behaviour, and the
reliability of empirical comparisons.

This aim is addressed through the following objectives.

\begin{description}[leftmargin=0pt,style=nextline]
    \item[O1.]
    To develop an explicit graph representation whose individual features jointly encode the local neighbourhoods around node pairs and the longer-range structure connecting them, while incorporating continuous node attributes without discretisation.
    \item[O2.]
    To evaluate whether such an explicit representation can provide accurate,
    scalable, and interpretable graph-level features across molecular,
    citation, and larger-scale graph benchmarks.

    \item[O3.]
    To design a graph-conditioned diffusion generator that uses graph-level
    structural information to guide the generation of discrete graphs.

    \item[O4.]
    To separate node-level diffusion, edge prediction, and constraint-aware
    decoding so that the graph generation process is inspectable and can
    incorporate feasibility constraints.

    \item[O5.]
    To develop a diagnostic evaluation framework that compares generative
    models with reliability checks and diagnoses multiple generative failure
    modes, including poor task fidelity,
    distributional mismatch, low augmentation value, local structural
    mismatch, and memorisation.

    \item[O6.]
    To compare generative models using uncertainty-aware summaries and
    reliability checks rather than relying only on single scalar metrics.
\end{description}

These objectives lead to the following research questions.

\begin{description}[leftmargin=0pt,style=nextline]
    \item[RQ1.]
    How can graph kernels be extended to incorporate continuous node attributes
    directly while retaining explicit, sparse, and interpretable feature maps?

    \item[RQ2.]
    How can graph-level structural information be used to condition a
    diffusion-based generator for labelled graphs?

    \item[RQ3.]
    How can graph generative models be evaluated in a way that distinguishes
    predictive utility, distributional alignment, memorisation, and structural
    fidelity?

    \item[RQ4.]
    How can empirical generator comparisons be made more reliable through
    uncertainty-aware summaries and statistical guarantees?
\end{description}

These questions are deliberately connected. The first question concerns the
representation of graph structure. The second uses such representations as
conditioning signals for generation. The third and fourth address the problem
of deciding whether generated samples should be trusted.

\begin{table}[t]
\centering
\footnotesize
\begin{tabularx}{\textwidth}{@{}l l l X@{}}
\toprule
Research question & Objectives & Chapter & Evidence \\
\midrule
RQ1 & O1--O2 & Chapter~\ref{ch:nsppk} &
Kernel construction, feature analysis, classification experiments, scalability
tests, and interpretability analysis. \\
RQ2 & O3--O4 & Chapter~\ref{ch:generator} &
Synthetic cycle--leg generation, conditional decoding diagnostics, ablations,
and QM9 molecular experiments. \\
RQ3 & O5 & Chapter~\ref{ch:rankgen} &
Quality, Utility, Indistinguishability, Similarity, and memorisation-sensitive
diagnostics across controlled perturbations and generative benchmarks. \\
RQ4 & O6 & Chapter~\ref{ch:rankgen} &
PAC-style reliability checks, repeated resampling, quantile summaries, and
model ranking analyses. \\
\bottomrule
\end{tabularx}
\caption{Mapping between research questions, objectives, thesis chapters, and
the main evidence used to answer each question.}
\label{tab:intro:rq-objective-map}
\end{table}

\section{Thesis Contributions}
\label{sec:intro:contributions}

This thesis makes three main methodological contributions.

\textbf{Explicit attributed graph representation.}
Chapter~\ref{ch:nsppk} introduces the Neighborhood Subgraph Pairwise Path
Kernel (NSPPK), an explicit graph kernel whose features jointly encode the
local neighbourhoods around node pairs and the broader graph structure
connecting them. Connector subgraphs are constructed from unions of shortest
paths, while continuous node attributes are incorporated directly into the
feature map without discretisation. The resulting representation is sparse,
deterministic, and interpretable, supports efficient dot-product kernel
evaluation, and allows predictive signals to be traced back to structural
graph patterns.

\textbf{Graph-conditioned Diffusion generation.}
Chapter~\ref{ch:generator} introduces the Graph-Conditioned Diffusion Generator
(GCDG), a conditional generator that separates continuous node-level denoising
from discrete graph construction. A graph-level condition vector guides a
transformer-based denoising model, while node existence, node and edge labels,
degrees, and edge probabilities are predicted through dedicated modules before
constraint-aware decoding. This staged design makes intermediate predictions
inspectable and allows graph feasibility constraints to be enforced explicitly
during generation.

\textbf{Statistically grounded generator evaluation.}
Chapter~\ref{ch:rankgen} introduces RankGen, a classifier-based framework for
evaluating generative models. RankGen defines complementary measures of
Quality, Utility, Indistinguishability, and Similarity, each designed to expose
a different failure mode. It combines these measures with PAC-style reliability
bounds, robust quantile summaries, and a ranking procedure that distinguishes
statistical confidence from fine-grained model comparison. The framework is
validated across controlled perturbations and generative tasks involving
molecular graphs, images, and open-domain text.

Together, these contributions support a qualified central claim: graph
generation benefits from treating representation, synthesis, and diagnosis as
connected design decisions. This is a methodological perspective rather than a
claim that the thesis evaluates one fully integrated system throughout. In the
reported generator experiments, graph-level descriptors are used to condition
GCDG; RankGen is developed as a broader diagnostic framework for evaluating
generated data rather than as an optimiser inside the generator. Generated
graphs should therefore not be judged using a single measure of plausibility,
but according to whether they preserve relevant structure, satisfy domain
constraints, support downstream use, and pass the reliability checks required
by the evaluation setting.

\section{Thesis Structure}
\label{sec:intro:structure}

Chapter~\ref{ch:theory} introduces the theoretical and methodological
background required for the thesis, including graph notation, graph
representation learning, graph kernels, neural graph models, diffusion-based
generation, and principles for evaluating generated data.

Chapter~\ref{ch:nsppk} develops NSPPK and presents its construction,
computational analysis, implementation, and empirical evaluation.

Chapter~\ref{ch:generator} develops GCDG and describes its conditioning
mechanism, corruption and diffusion processes, prediction modules, decoding
procedure, training objective, and empirical evaluation.

Chapter~\ref{ch:rankgen} develops RankGen, defines its evaluation measures and
reliability guarantees, and assesses the framework on molecular graphs, images,
and open-domain text, as well as under controlled perturbation settings.

Chapter~\ref{ch:conclusion} summarises the findings, discusses the limitations
of the proposed methods, and outlines directions for future research.
